% ============================================================
\chapter{Applications}
\label{ch:applications}
% ============================================================

\begin{intuition}[title=\textbf{Central idea}]
Bayesian statistics is ubiquitous in modern applications: clinical trials,
A/B testing, automatic hyperparameter optimization, Bayesian neural networks,
and many scientific domains. This chapter illustrates the power of the
Bayesian paradigm on concrete problems.
\end{intuition}

% ------------------------------------------------------------
\section{Bayesian clinical trials}
% ------------------------------------------------------------

\begin{definition}[Adaptive Bayesian clinical trial]
\index{Bayesian clinical trial}
An \textbf{adaptive Bayesian clinical trial} is a trial in which the design
(sample size, patient allocation, stopping criteria) is modified during the
study based on the posterior distribution of the efficacy parameters.
\end{definition}

\begin{example}
Let $\theta_T$ be the response probability under treatment and $\theta_C$
under control. We model:
\[
  \theta_T \sim \mathrm{Beta}(1, 1), \quad
  \theta_C \sim \mathrm{Beta}(1, 1).
\]
After observing $s_T$ successes out of $n_T$ treated patients and $s_C$
out of $n_C$ controls:
\[
  \theta_T \mid \text{data} \sim \mathrm{Beta}(1 + s_T,\, 1 + n_T - s_T).
\]
We stop for efficacy if $\PP(\theta_T > \theta_C \mid \text{data}) > 0.99$,
or for futility if $\PP(\theta_T > \theta_C + \delta \mid \text{data}) < 0.05$.
\end{example}

\begin{remark}
Adaptive Bayesian trials have been approved by the FDA since the 2010s.
They reduce the required sample size by 20--40\% on average compared to
fixed designs.
\end{remark}

\begin{minted}{python}
import numpy as np
from scipy.stats import beta

def prob_superiority(s_T, n_T, s_C, n_C, n_samples=100000):
    """P(theta_T > theta_C | data) via Monte Carlo."""
    theta_T = beta.rvs(1 + s_T, 1 + n_T - s_T, size=n_samples)
    theta_C = beta.rvs(1 + s_C, 1 + n_C - s_C, size=n_samples)
    return np.mean(theta_T > theta_C)
\end{minted}

% ------------------------------------------------------------
\section{Bayesian A/B testing}
% ------------------------------------------------------------

\begin{definition}[Bayesian A/B test]
\index{A/B testing}
A \textbf{Bayesian A/B test} compares two variants (A and B) of a product
by maintaining a posterior distribution on each variant's conversion rate.
A decision is made when the probability that one variant is better exceeds
a fixed threshold.
\end{definition}

\begin{proposition}[Expected loss]
\index{expected loss}
The \textbf{expected loss} from choosing B is:
\[
  \mathcal{L}_B = \E\!\bigl[\max(\theta_A - \theta_B, 0) \mid \text{data}\bigr].
\]
We choose B if $\mathcal{L}_B < \epsilon$ (acceptable loss threshold).
This criterion directly controls the cost of a wrong decision.
\end{proposition}

\begin{attention}
Unlike the frequentist test (p-value), the Bayesian A/B test allows
``peeking'' (examining data during collection) without inflating the error
rate, because the posterior is coherent at any time.
\end{attention}

% ------------------------------------------------------------
\section{Bayesian optimization}
% ------------------------------------------------------------

\begin{definition}[Bayesian optimization]
\index{Bayesian optimization}
\textbf{Bayesian optimization} seeks the global minimum of an expensive
function $f : \mathcal{X} \to \R$ by building a probabilistic surrogate
model (typically a GP) and selecting evaluation points via an
\textbf{acquisition function}.
\end{definition}

\begin{theorem}[GP posterior as surrogate]
Let $f \sim \mathcal{GP}(0, k)$ and $\mathcal{D}_n = \{(x_i, y_i)\}_{i=1}^n$
with $y_i = f(x_i) + \epsilon_i$. The prediction at $x$ is:
\begin{align}
  \mu_n(x) &= \mathbf{k}(x)^\top (\mathbf{K} + \sigma^2 \mathbf{I})^{-1}\mathbf{y}, \\
  \sigma_n^2(x) &= k(x,x) - \mathbf{k}(x)^\top (\mathbf{K} + \sigma^2 \mathbf{I})^{-1}\mathbf{k}(x).
\end{align}
\end{theorem}

\begin{definition}[Acquisition functions]
Common acquisition functions include:
\begin{itemize}
  \item \textbf{Expected Improvement (EI)}:
    $\mathrm{EI}(x) = \E\!\bigl[\max(f_{\min} - f(x), 0)\bigr]$.
  \item \textbf{Upper Confidence Bound (UCB)}:
    $\mathrm{UCB}(x) = \mu_n(x) - \kappa\,\sigma_n(x)$ (for minimization).
  \item \textbf{Probability of Improvement (PI)}:
    $\mathrm{PI}(x) = \Phi\!\left(\frac{f_{\min} - \mu_n(x)}{\sigma_n(x)}\right)$.
\end{itemize}
\end{definition}

\begin{example}
Hyperparameter optimization for a neural network:
\begin{minted}{python}
from skopt import gp_minimize

def objective(params):
    lr, dropout = params
    model = build_model(lr=lr, dropout=dropout)
    return -model.fit_and_evaluate(X_train, y_train, X_val, y_val)

result = gp_minimize(
    objective,
    dimensions=[(1e-5, 1e-1, "log-uniform"),  # learning rate
                (0.0, 0.5)],                    # dropout
    n_calls=50,
    random_state=42
)
\end{minted}
\end{example}

% ------------------------------------------------------------
\section{Bayesian neural networks}
% ------------------------------------------------------------

\begin{definition}[Bayesian neural network (BNN)]
\index{Bayesian neural network}
A \textbf{BNN} places a prior $p(\mathbf{w})$ on the network weights
$\mathbf{w}$ and computes the posterior $p(\mathbf{w} \mid \mathcal{D})$.
The prediction is:
\[
  p(y_* \mid x_*, \mathcal{D}) = \int p(y_* \mid x_*, \mathbf{w})\,
  p(\mathbf{w} \mid \mathcal{D})\,d\mathbf{w}.
\]
\end{definition}

\begin{remark}
The above integral is intractable for deep networks. Three practical
approaches:
\begin{enumerate}
  \item \textbf{MC Dropout}: use dropout at inference time as a variational
        approximation (Gal and Ghahramani, 2016).
  \item \textbf{Bayes by Backprop}: variational inference with
        reparameterization on each weight.
  \item \textbf{Deep ensembles}: train $M$ networks and average
        (frequentist approximation of uncertainty).
\end{enumerate}
\end{remark}

\begin{proposition}[Epistemic vs.\ aleatoric uncertainty]
The predictive variance of a BNN decomposes as:
\[
  \underbrace{\Var_{p(\mathbf{w} \mid \mathcal{D})}\!\bigl[\E[y_* \mid x_*, \mathbf{w}]\bigr]}_{\text{epistemic}}
  + \underbrace{\E_{p(\mathbf{w} \mid \mathcal{D})}\!\bigl[\Var[y_* \mid x_*, \mathbf{w}]\bigr]}_{\text{aleatoric}}.
\]
Epistemic uncertainty decreases with more data; aleatoric does not.
\end{proposition}

% ------------------------------------------------------------
\section{Applications in astronomy}
% ------------------------------------------------------------

\begin{example}
In cosmology, estimating the parameters of the $\Lambda$CDM model
(Hubble constant $H_0$, matter density $\Omega_m$, etc.) relies on
Bayesian inference applied to cosmic microwave background (CMB) data.
The software \texttt{CosmoMC} uses MCMC to explore the posterior in
$\sim 30$ dimensions.
\end{example}

\begin{example}
Exoplanet detection via the radial velocity method models the signal as:
\[
  v(t) = \sum_{j=1}^{K} K_j \sin(2\pi t / P_j + \phi_j) + \epsilon(t),
\]
where the number of planets $K$ is unknown. Bayesian model comparison
(Bayes factors) determines $K$.
\end{example}

% ------------------------------------------------------------
\section{Applications in ecology}
% ------------------------------------------------------------

\begin{example}
Bayesian \textbf{capture-recapture} models estimate population size $N$.
With a prior $N \sim \mathrm{Poisson}(\lambda)$ and recapture data, the
posterior on $N$ naturally quantifies uncertainty---essential for
conservation management.
\end{example}

\begin{definition}[Hierarchical model in ecology]
\index{hierarchical model}
\textbf{Bayesian hierarchical models} are standard in ecology for modeling:
\begin{itemize}
  \item between-site variability (spatial random effects),
  \item imperfect detection (occupancy models),
  \item temporal dynamics (state-space models).
\end{itemize}
The parameters for each site $j$ share a common hyperprior:
$\theta_j \sim p(\theta_j \mid \psi)$, $\psi \sim p(\psi)$.
\end{definition}

% ------------------------------------------------------------
\section{Key formulas}
% ------------------------------------------------------------

\begin{keyformulas}
\begin{align}
  \text{EI:} \quad & \mathrm{EI}(x) = (f_{\min} - \mu_n(x))\,\Phi(z) + \sigma_n(x)\,\phi(z),
    \quad z = \frac{f_{\min} - \mu_n(x)}{\sigma_n(x)} \\
  \text{BNN prediction:} \quad & p(y_* \mid x_*) \approx
    \frac{1}{M}\sum_{m=1}^M p(y_* \mid x_*, \mathbf{w}^{(m)}),
    \quad \mathbf{w}^{(m)} \sim p(\mathbf{w} \mid \mathcal{D}) \\
  \text{A/B loss:} \quad & \mathcal{L}_B = \E[\max(\theta_A - \theta_B, 0) \mid \text{data}]
\end{align}
\end{keyformulas}

% ------------------------------------------------------------
\section{Exercises}
% ------------------------------------------------------------

\begin{exercise}
Implement a Bayesian A/B test with Beta(1,1) priors to compare click-through
rates of two web pages. Simulate data and plot the evolution of
$\PP(\theta_A > \theta_B \mid \text{data})$ across observations.
\end{exercise}

\begin{exercise}
Perform Bayesian optimization to find the minimum of the Branin function
$f(x_1, x_2)$ using a GP with Mat\'ern kernel. Compare with random search
in terms of the number of evaluations needed.
\end{exercise}

\begin{exercise}
Implement MC Dropout on a two-hidden-layer neural network for a regression
problem. Plot the mean prediction and the 95\% confidence interval. Observe
how uncertainty behaves far from training data.
\end{exercise}

\begin{exercise}
Design an adaptive Bayesian clinical trial with the following rules: stop
for efficacy if $\PP(\theta_T - \theta_C > 0.05 \mid \text{data}) > 0.95$,
stop for futility if $\PP(\theta_T > \theta_C \mid \text{data}) < 0.2$.
Simulate 1000 trials under $H_0$ and $H_1$ and estimate the operating
characteristics (power, type I error rate, average sample size).
\end{exercise}

\begin{exercise}
In an exoplanet detection problem, compare models $M_0$ (pure noise) and
$M_1$ (one planet). Compute the Bayes factor $\mathrm{BF}_{10}$ by Monte
Carlo integration on simulated data. Discuss Jeffreys' scale for
interpretation.
\end{exercise}
